\chapter[p-adic Fields]{$p$-adic Fields}

\section[p-adic Equations]{$p$-adic Equations}\label{sec:padic_equations}

\begin{definition}\label{def:primitive}
\lean{Function.IsPrimitive}
\leanok
Let $R$ be a ring. An element $(x_1, \cdots, x_n) \in R^n$ is \emph{primitive} if at least one of the $x_i$ is invertible.
\end{definition}


\begin{proposition}\label{prop:common_root}
\lean{Function.IsPrimitive}
\leanok
\uses{def:primitive}
Let $f^{(i)} \in \Z_p[X_1, \cdots, X_m]$ be homogeneous polynomials with $p$-adic integer coefficients. For each $n \ge 1$, denote by $f_n^{(i)}$ the reduction of $f^{(i)}$ modulo $p^n$.
 The following are equivalent:
\begin{enumerate}[(a)]
\lean{Padic.common_root_tfae}
\leanok
\item The $f^{(i)}$ have a nontrivial common zero in $\Q_p^m$.
\item The $f^{(i)}$ have a common primitive zero in $\Z_p^m$.
\item For all $n \ge 1$, the $f_n^{(i)}$ have a common primitive zero in $(\Z / p^n \Z_p)^m$.
\end{enumerate}
\end{proposition}

\begin{proof}
The implication (b) $ \implies$ (a) is trivial. Conversely, if $x = (x_1, \cdots, x_m) \in \Q_p^m$ is a nontrivial common zero of the $f^{(i)}$, put
\[ h := \min(v_p(x_1), \cdots, v_p(x_m)) \quad \text{and} \quad y := p^{-h} x. \]
Then $y \in \Z_p^m$ is a common primitive zero of the $f^{(i)}$.
For the equivalence (b) $\iff$ (c), Serre applies the first Lemma from \cite[\S II.2.1]{Ser73}. However, using \mathlib's \texttt{DirectLimit} definition here might be tricky. It could be easier to just prove this specific result; note that the file \texttt{Mathlib.NumberTheory.Padics.RingHoms} contains useful results relating $\Z_p$ and $\Z/p^n\Z$ (without relying on \texttt{DirectLimit}).

\end{proof}

\section[Squares in the p-adic units]{Squares in $\Q_p^{\times}$}\label{sec:padic_squares}

TODO: expand this section.

\begin{theorem}\label{thm:padic_squares_open}
\lean{Padic.squares}
\leanok
${\Q_p^\times}^2$ is an open subgroup of $\Q_p^\times$.
\end{theorem}

\section[Units and Legendre symbols]{Units and Legendre symbols}\label{sec:padic_legendre}

\begin{definition}\label{def:padicUnitPart_qp}
\lean{Padic.padicUnitPart_qp}
\leanok
Let $p$ be a prime and $x \in \Q_p^\times$. Let $v$ denote the usual $p$-adic valuation.
The \emph{unit part in $\Q_p$} of $x$ is defined as
\[x \cdot p^{-v_p(x)}.\]
This is an element in $\Q_p^\times$.
\end{definition}

\begin{lemma}\label{lem:padic_unit_part_unit}
\lean{Padic.padic_unit_part_unit}
\uses{def:padicUnitPart_qp}
\leanok
Let $p$ be a prime and $x \in \Q_p^\times$. Then $\|x \cdot p^{-v_p(x)}\| = 1$, so the
unit part is a unit, i.e. it is in $(\Z_p)^\times$.
\end{lemma}
\begin{proof}
Straightforward.
\end{proof}

\begin{definition}\label{def:padicUnitPart}
  \lean{Padic.padicUnitPart}
  \uses{def:padicUnitPart_qp,lem:padic_unit_part_unit}
  \leanok
Let $p$ be a prime and $x \in \Q_p^\times$. The \emph{unit part} of $x$ is defined as the element
of $\Z_p^\times$ corresponding to $x \cdot p^{-v_p(x)}$.
\end{definition}

% \begin{lemma}
% 	\label{lem:is_twoadic_unit_odd_p}
% 	\lean{Padic.is_twoadic_unit_odd_p}
% 	\leanok
% Let $p$ be an odd prime. Then its image in $\Z_2$ is a unit.
% \end{lemma}

% \begin{proof}
% Immediate, since $ 2 \nmid p$.
% \end{proof}

% \begin{definition}
% 	\label{def:p2}
% 	\lean{Padic.p2}
% 	\uses{lem:is_twoadic_unit_odd_p}
% 	\leanok
% Let $p$ be an odd prime. Then we define $p_{in_Q2} \in (\Z_2)^\times$ as the image of $p$ in
% $(\Z_2)^\times$.
% \end{definition}

\begin{definition}\label{def:legendreSym'}
  \lean{Padic.legendreSym'}
  \leanok
Let $p$ be a prime and $u \in \Z_p^\times$. The \emph{Legendre symbol} $\left(\dfrac{u}{p}\right)$
is defined as
\[\left(\dfrac{u}{p}\right)=\left(\dfrac{\overline{u}}{p}\right),\]
where $\overline{u}$ is the image of $u$ under the natural map $\Z_p^\times \to (\Z/p\Z)^\times$.
\end{definition}
